\subsection{Versuchsaufbau}

Das Paper von Feldman und Golumbic \cite{feldmangolumbic} wählte einen \textit{brute force}-Algorithmus als Baseline für Performance-Vergleiche.
Da ein \textit{brute force} Algorithmus bei der großen Anzahl der Vorlesungen im VVZ der WU eine zu lange Laufzeit hätte, wurde eine \textit{Integer Linear Programming} (ILP) Lösung als neue Baseline gewählt. Diese gewährleistet, ebenso wie der \textit{brute force}-Ansatz für kleine Instanzen, die Ermittlung des optimalen Stundenplans (mit der höchsten \textit{Mark}). Die Ergebnisse der heuristischen Algorithmen (\textit{Hill-Climbing v1}, \textit{Hill-Climbing v3}, \textit{Offering-Order}) werden folglich relativ zur optimalen Lösung des ILP-Ansatzes bewertet. Zusätzlich zu den heuristischen Algorithmen wird auch eine GPU-optimierte Variante des ILP-Ansatzes getestet.

Als Maß für die Schwierigkeit (\textit{Difficulty}) der Probleminstanz wird in dieser Arbeit nicht die Laufzeit des \textit{brute force}-Algorithmus verwendet, sondern die Anzahl der zu planenden Kurse (\textit{N}).

Die Algorithmen wurden anhand der folgenden Leistungskennzahlen evaluiert:
\begin{enumerate}
    \item \textbf{Mark}: Die Qualität der gefundenen Lösung (in Prozent relativ zum Baseline Ansatz).
    \item \textbf{Laufzeit}: Die zur Lösungsfindung benötigte Zeit.
    \item \textbf{Speicherbedarf}: Der benötigte Hauptspeicher.
\end{enumerate}


\subsubsection{Zielfunktion}

Die \textit{Mark} dient als Bewertungsfunktion für die Qualität eines erstellten Stundenplans und entspricht der im Paper von Feldman und Golumbic \cite{feldmangolumbic} definierten Zielfunktion. Sie quantifiziert, wie gut ein Stundenplan die Präferenzen und \textit{Constraints} der Studierenden erfüllt. Die Berechnung erfolgt ausschließlich für \textbf{gültige} Stundenpläne, d.\,h. solche, die keine restriktiven \textit{Constraints} (\(|p| = P\)) verletzen.

Jede Präferenz oder Restriktion ist im Modell durch eine \textit{Priorität} \(p\) im Wertebereich \([-P, P]\) definiert. Positive Prioritäten kennzeichnen wünschenswerte Eigenschaften (z.\,B. bevorzugte Lehrveranstaltungen oder Zeitfenster), negative Prioritäten hingegen unerwünschte. Die Stärke des Präferenzwerts ist proportional zum Absolutwert der Priorität.

Die \textit{Mark} eines gültigen Stundenplans \(S\) ergibt sich aus der Summe der positiven Beiträge durch gewählte Kurse und den Abzügen für Verletzungen sogenannter \textit{fixed} und \textit{non-fixed Constraints}:

\begin{equation}
    Mark(S) =
    \underbrace{\sum_{c \in C_S} q_c}_{\text{Kursprioritäten}}
    - \underbrace{\sum_{h \in H_{\text{violated}}^{(f)}} |p_h^{(f)}|}_{\text{Strafen für fixed-Constraints}}
    - \underbrace{\sum_{t \in T_{\text{violated}}^{(n)}} v_t \cdot |p_t^{(n)}|}_{\text{Strafen für non-fixed-Constraints}},
    \label{eq:mark}
\end{equation}

wobei gilt:

\begin{itemize}
    \item \(C_S\): Menge aller im Stundenplan \(S\) enthaltenen Kurse.
    \item \(q_c\): Priorität des Kurses \(c\).
    \item \(H_{\text{violated}}^{(f)}\): Menge aller Stunden, in denen ein \textit{fixed Constraint} verletzt wurde.
    \item \(p_h^{(f)}\): Priorität der Stunde \(h\). Eine Verletzung liegt vor, wenn
          \begin{itemize}
              \item \(p_h^{(f)} < 0\) und die Stunde aktiv ist (der Studierende hat in einer ungewünschten Stunde Unterricht), oder
              \item \(p_h^{(f)} > 0\) und die Stunde inaktiv ist (der Studierende hat in einer gewünschten Stunde keinen Unterricht).
          \end{itemize}
    \item \(T_{\text{violated}}^{(n)}\): Menge der verletzten \textit{non-fixed constraints}.
    \item \(p_t^{(n)}\): Priorität des \textit{Constraint} \(t\).
    \item \(v_t\): Anzahl der verletzten Stunden (bzw. des Zeitumfangs), der durch \textit{Constraint} \(t\) betroffen ist.
\end{itemize}

Je höher der Wert von \(M(S)\), desto besser erfüllt der Stundenplan die angegebenen Präferenzen. Ziel der Optimierungsalgorithmen ist daher die Maximierung von \(M(S)\).


\subsubsection{Durchführung}

Die experimentelle Evaluation der Algorithmen erfolgte durch systematische Benchmarks, die Performance und Lösungsqualität der implementierten Ansätze unter verschiedenen Szenarien und Schwierigkeitsgraden testen.

Die Experimente wurden auf dem Hybridair-GPU-Server der Wirtschaftsuniversität Wien durchgeführt. Die Hardware Infrastruktur hat folgende Spezifikationen:

\begin{itemize}
    \item \textbf{Prozessor:} AMD EPYC 9334
    \item \textbf{Arbeitsspeicher:} 128 GB RAM
    \item \textbf{Grafikkarte (ILP GPU):} Zwei NVIDIA RTX A5000 GPUs (jeweils 24 GB VRAM) mit CUDA 12.4
    \item \textbf{Betriebssystem:} Ubuntu/Debian via Proxmox Kernel (6.8.12-11-pve)
\end{itemize}

Um eine verlässliche und gleichbleibende Testumgebung zu garantieren, wurde das Tool \textbf{benchexec} genutzt. Jeder Testlauf hatte ein Zeitlimit von 10 Sekunden und lief auf einem Prozessorkern. Um die Ergebnisse nicht zu verfälschen, wurde der Server während der Tests nicht anderweitig verwendet.

Die Algorithmen wurden in 14 verschiedenen Szenarien getestet. Jedes Szenario stellt eine vorkonfigurierte Zusammensetzung von \textit{Constraints} dar (definiert in JSON-Konfigurationsdateien, z.B. \texttt{constraint1.json}).

Der Schwierigkeitsgrad (\textit{Difficulty}) einer Probleminstanz wird durch die Anzahl der zu planenden Kurse $N$ definiert, im Gegensatz zur Laufzeit des \textit{brute force}-Algorithmus, die in der Originalarbeit verwendet wurde. Die Tests wurden inkrementell für jede mögliche Kursanzahl $N$, beginnend bei $N=1$ bis zur maximal möglichen Kursanzahl im jeweiligen Szenario, durchgeführt. Die maximale Kursanzahl wurde dabei vorab mithilfe des ILP-Baseline-Ansatzes ermittelt.

% \paragraph{Ablauf der Messung}

Für jeden Algorithmus ($A \in \{ \text{ILP}, \allowbreak \text{ILP GPU}, \allowbreak \text{Offering-Order}, \allowbreak \text{Hill-Climbing v1}, \allowbreak \text{Hill-Climbing v3} \}$), für jede Kursanzahl $N$ und für jedes Szenario $S$ wurde die Messreihe wie folgt durchgeführt:

\begin{enumerate}
    \item Jeder Algorithmus wurde für eine feste Kursanzahl $N$ 50 Mal unabhängig voneinander ausgeführt.
    \item Die 50 Messungen pro Algorithmus und Schwierigkeitsgrad wurden sequenziell durchgeführt, um sicherzustellen, dass die Messergebnisse für die Laufzeit und den Speicherbedarf unabhängig sind.
    \item Es wurden folgende Metriken für jede der 50 Ausführungen erfasst:
          \begin{itemize}
              \item Laufzeit ($\mathit{t}$): Die zur Lösungsfindung benötigte Zeit (in Sekunden).
              \item Hauptspeicher ($\mathit{M}$): Der während der Ausführung belegte Hauptspeicher (in Megabyte).
              \item Qualität ($\mathit{Mark}$)
          \end{itemize}
    \item Für die 50 Messwerte der Laufzeit und des Hauptspeichers wurden folgende statistische Kennzahlen berechnet:
          \begin{itemize}
              \item Median ($t_{\text{median}}, M_{\text{median}}$)
              \item Mittelwert ($t_{\text{mean}}, M_{\text{mean}}$)
              \item Minimalwert ($t_{\text{min}}, M_{\text{min}}$)
              \item Maximalwert ($t_{\text{max}}, M_{\text{max}}$)
              \item Standardabweichung ($t_{\text{stdev}}, M_{\text{stdev}}$)
          \end{itemize}
\end{enumerate}

% \paragraph{Abgrenzung der Messung}

Um die Algorithmen vergleichbar zu halten, wurde der gemessene Zeit- und Speicherbedarf auf die Ausführung des Algorithmus selbst begrenzt. Das Preprocessing — insbesondere die Vorabberechnung der \textit{Mark} für jedes Offering, die einmalig vor der Hauptschleife durchgeführt wird — ist nicht in den Messwerten enthalten.

